%% Copyright 2026
%% Licensed under Creative Commons Attribution–ShareAlike 4.0 International (CC BY-SA 4)

\documentclass[12pt]{article}
\usepackage{titling}
\usepackage{listings}
\usepackage{xcolor}
%% \usepackage{fontspec}

\usepackage{sourcecodepro}

\newcommand{\NULL}{{\texttt{NULL}}}%
\newcommand{\cvolatile}{{\texttt{volatile}}}%
\newcommand{\cdefine}{{\texttt{\#define}}}%
\newcommand{\cstatic}{{\texttt{static}}}%
\newcommand{\cconst}{{\texttt{const}}}%
\newcommand{\cCC}{{\texttt{C/C++}}}%
\newcommand{\cC}{{\texttt{C}}}%

\lstset{ %
   %%numbers=left,   
   %%stepnumber=2,
   %%numbersep=5pt,
   %% numberstyle=\footnotesize,
   language=C++,
   basicstyle=\footnotesize\ttfamily,
   breaklines=true,
   frame=shadowbox,
}

\author{spinoff@protonmail.com}
\date{August 27, 2021}

\title{So, you want to be an \\embedded software engineer?}

\begin{document}
\begin{titlingpage}
\maketitle
\begin{abstract}
This document focuses on the core
competency of C language programming and basics for 
embedded software engineering.  Among many other things
that are important, strong knowledge of C is critical for
success.
\end{abstract}
\par
\vskip1cm
{\textsc{version 1.1}}
\vskip5mm

\end{titlingpage}

\section{Introduction}

Where noted ($\ddagger$), these questions were stolen from an article
in Embedded Systems Journal in 2002.

When I first read the questions back then it occurred to me that some work needed to be
done within the company I was with at the time to enhance our test for prospective 
engineers.  I've dragged this list around with me since.

The nature of these questions is not to be misinterpreted -- a candidate should know
how to solve all these problems. Otherwise, they usually are not hired.

Here we go! 

\section{The C-Test}

\begin{enumerate}

\item $\ddagger$ Using the \cdefine\ directive, how would you declare a manifest
constant that is the number of seconds in a year (Disregard leap year in your
answer).
\par

\item $\ddagger$ Write the ``standard'' {\texttt{MIN}} macro such that the macro takes two arguments
and resolves to the value of the minimum of the arguments.
\newpage
Comment on this code:
\begin{lstlisting}
int a;
int b;
int *pa = &a;

// ...

int r = MIN(*pa++, b);

\end{lstlisting}

\item $\ddagger$ What is the purpose of the pre-processor directive {\texttt{\#error}} ?


\item $\ddagger$ Infinite loops often arise in embedded systems. How do you code an infinite loop?


\item $\ddagger$ Using the symbol {\texttt{x}}, give the declarations for the following:
\begin{itemize}
\item An integer.
\item A pointer to an integer.
\item A pointer to a pointer to an integer.
\item An array of ten integers.
\item An array of ten pointers to integers.
\item A pointer to an array of ten integers
\item A pointer to a function that takes an integer argument  and returns an integer.
\item An array of ten pointers to functions that each take an integer argument and each return an integer.
\end{itemize}

\item $\ddagger$ What are the uses of the keyword {\texttt{static}} ?


\item $\ddagger$ What does the keyword \cconst\  mean?

Also, this leads to a set of questions here that should be addressed.
Comment on what each of these declarations mean:

\begin{itemize}
\item  {\texttt{const int x;}}
\item {\texttt{int const x;}}
\item {\texttt{int* const x;}}
\item {\texttt{const int* x;}}
\item {\texttt{const int* const x;}}
\end{itemize}

\end{enumerate}

\section{Getting Warmed Up}

By now you should have a good idea of the candidate's {\texttt{C}} capabilities.
But we still need to dig deeper.
\par
\begin{enumerate}
\item $\ddagger$ What does the keyword {\texttt{volatile}} mean?  Give three examples
of its use.

Can a ``variable'' in C be both {\texttt{volatile}}
and {\texttt{const}}?
Ask the candidate to comment on this code:
\begin{lstlisting}
const volatile int* preg;
\end{lstlisting}

\item $\ddagger$ Embedded systems often require the manipulation of bits in
registers or other parameters.  Given a variable 
{\texttt{unsigned int x}}, show how to:
  \begin{itemize}
    \item Set bit 3 of {\texttt{x}} (the fourth bit)
    \item Clear bit 3 of {\texttt{x}}
  \end{itemize}

\item $\ddagger$ Comment on this code and explain what is the output (and why)?

\begin{lstlisting}
void foo(void)
{
    unsigned int a = 6;
    int b = -20;
    
    ( a + b  > 6 ) ? puts("> 6\n") :
                     puts("<= 6\n");
}
\end{lstlisting}

\end{enumerate}

\section{Downhill now}

At this point we are nearing the end of the test.  If the candidate
is demoralized then that's a bad sign. They should have been able
to whip through these questions so far without a sweat.

But now we're going to get more elaborate in our questions and stray
beyond the pure syntax and use of {\texttt{C}} in the embedded system
scenario.

\begin{enumerate}
\item $\ddagger$ {\texttt{typedef}} is frequently used in C to declare a complex
data structure.   It's also possible to use the pre-processor to do
something similar.  For instance, consider the following code:

\begin{lstlisting}
#define dPS struct s *

// or..

typedef struct s* tPS;
\end{lstlisting}

\item Although it is not common in an embedded system, there are RTOS
in the industry that do synthesize some dynamic memory allocation
routines.  So, explain the the problem of using dynamic memory
in an embedded software application.  Describe the issues. Start 
with this simple question -- explain the difference between heap, 
free-store, and the stack?

A follow up question is:   Can you explain why I would want to avoid
too many Tasks in an application that uses an RTOS?  How many is too many?
What else should I be concerned about when I have a lot of tasks in an
embedded software application?

\item What is the difference between Heap and Stack?

\item Comment on this code.  What does it do?  Why?

\begin{lstlisting}
char* ptr;
if (( ptr = (char *) malloc (0)) == NULL )
{
     puts ( " got a null pointer " );
{
else
{
     puts ( " got a valid pointer " );
}
\end{lstlisting} 

\end{enumerate}

%%
%%
%%\section{Abstraction}
%%
%%In this section we're going to cover more abstract issues and concepts
%%that can be asked during the interview for a potential embedded software
%%engineer (or any software engineer for that matter).
%%
%%These questions are not specific to embedded software, but there is
%%a common set of knowledge that software engineers should know before
%%moving forward in the review process.
%%
%%\par
%%
%%\begin{enumerate}
%%
%%\item What makes a good API?
%%\item What features in C++ make it object-oriented?
%%\item How is data hidden from other classes in C++ ?
%%\item What is a virtual function in C++?
%%\item (C++ only) What is the difference between a pointer and a reference?
%%\item What are the key differences between C and C++?
%%
%%%% Section
%%\end{enumerate} 

\end{document}

